\section{Encoded paths and finite local displacement procedures}
\label{sec:geometry}

The encoded-path and displacement-field method follows the construction
in Rubinstein~\cite[Sections~3.1--3.2]{Rubinstein2016Settling}: current and successor
blocks trace a copy--toggle polygon, and radial interpolation moves from
the tangent through the inward and reverse-tangent directions to a fixed
vertical direction. Here we prove the geometric interface at the normalized
codeword scale, including rounded corners, displacement after clipping,
and finite local access with average errors. This section is independent
of both circuit realizations. Its hypotheses concern entire binary words and local access
to their bits; they do not depend on the algebraic interpretation of an
individual coordinate. In particular, in the inner application the graph
is the path for the fixed true input tuple, and the distinguished source
is accessed through the virtual procedures of Section~\ref{sec:inner-gadget}.

\subsection{The geometric statement and its parameters}
\label{geo:setup}

Let $G$ be a finite directed graph of indegree and outdegree at most one,
with no self-edge and no isolated vertex. Fix a source $w_*$ with an actual
successor. All other vertices missing an incoming or an outgoing edge are
called desired endpoints. Suppose that distinct vertices have encodings
$E_w\in\{0,1\}^m$ satisfying
\begin{equation}\label{geo:distance}
 \norm{E_u-E_v}_{\rm b}\ge d_*>0\qquad(u\ne v).
\end{equation}
We decrease $d_*$ to a positive rational at most one. On any number of
blocks of length $m$ use the inner product and norm
\begin{equation}\label{geo:norm}
 \langle U,V\rangle_{\rm b}=m^{-1}\sum_j U_jV_j,
 \qquad \norm{U}_{\rm b}^2=m^{-1}\sum_j U_j^2.
\end{equation}
An unqualified norm in this section means this norm. The infinity norm
is unnormalized. Write $\mathcal B=[-1,2]^{4m}$, and let
$\clip_{\mathcal B}$ denote coordinatewise clipping to this cube.
Clipping is nonexpansive for~\eqref{geo:norm}. Scalars occupying a block
mean constant repetitions.

\begin{theorem}[Encoded-path geometry]\label{thm:geometry}
Under these hypotheses, there is a union $\mathcal L\subset\mathcal B$
of directed, rounded paths and cycles, with desired endpoint set
$\mathcal E$, having the following properties. For every sufficiently
small positive rational $\tau$, there are a positive integer $B_1$ and
positive rational constants $c_*,\beta_0$, depending only on $d_*$ and
$\tau$, and a field $v:\mathcal B\to\mathbb R^{4m}$ such that
\begin{align}
 \norm{v(X)}_\infty&\le B_1-1,\label{geo:bounded-field}\\
 \norm{\clip_{\mathcal B}(X+\beta v(X))-X}
   &\ge c_*\beta
 \quad\text{if }\dist(X,\mathcal E)>10\tau,
       \quad 0<\beta\le\beta_0.\label{geo:displacement}
\end{align}
The field is continuous on the relative open set
$\{X\in\mathcal B:\dist(X,\mathcal E)>10\tau\}$.
Every $P\in\mathcal E$ has first block exactly $E_w$ for a desired
endpoint $w$. Consequently, if $\norm{X-P}\le10\tau$, rounding $X_1$
at $1/2$ gives relative Hamming distance at most $4(10\tau)^2$ from
$E_w$. The exceptional top of the distinguished vertical is excluded
from $\mathcal E$ and has field $-\mathbf e$, where
$\mathbf e=(0,0,0,1)$.

The initial choice of $\tau$ may impose any further fixed positive upper
bounds, including validation and global-recovery radii. All these choices,
and $B_1,c_*,\beta_0$, precede every numerical accuracy choice and the
choice of a step size $\beta$.
\end{theorem}

We prove the theorem before formulating the precise access contract for
its finite numerical implementation. In particular no global regularity
claim at desired endpoints will be needed.

\subsection{Polygons, fillets, and local projection}
\label{geo:curves}

For an edge $u\to v$, draw, with fourth block zero,
\begin{equation}\label{geo:polygon}
 \begin{aligned}
 &(E_u,E_u,0)\xrightarrow{\;2\;}(E_u,E_v,0)
                 \xrightarrow{\;3\uparrow\;}(E_u,E_v,1)\\
 &\qquad\xrightarrow{\;1\;}(E_v,E_v,1)
                 \xrightarrow{\;3\downarrow\;}(E_v,E_v,0).
 \end{aligned}
\end{equation}
An arrow label names the changing block; the up and down arrows record
the direction of the third block. This is a schematic of block changes,
not a distance-preserving planar drawing.
Concatenate these polygons at graph-vertex junctions. Put
$H_*=(E_{w_*},E_{w_*},0)$ and prepend the segment
$(H_*,2)\longrightarrow(H_*,0)$ to the path out of $w_*$.
We call the four segment types in~\eqref{geo:polygon} \emph{horizontal},
meaning that their fourth block is zero; this includes the two third-block
toggles. Every horizontal segment has length in $[d_*,1]$, and consecutive unit
directions are orthogonal: the successive varying blocks in a polygon
are $2,3,1,3$, a graph junction changes from block $3$ to block $2$,
and the special bottom changes from block $4$ to block $2$.

\begin{lemma}[Separation of original segments]\label{geo:separation}
There is $d_{\rm sep}>0$, depending only on $d_*$, such that distinct
nonconsecutive original segments are at distance at least $d_{\rm sep}$.
Their adjacency graph, with adjacency meaning consecutiveness, has no
triangle.
\end{lemma}
\begin{proof}
For binary words $a,b,c$ and $0\le s\le1$,
\begin{equation}\label{geo:binary-chord}
 \norm{(1-s)a+sb-c}\ge\min(s,1-s)\norm{a-b}.
\end{equation}
Indeed on a coordinate where $a,b$ differ, the chord coordinate is $s$
or $1-s$, whose distance from a binary coordinate is at least
$\min(s,1-s)$. Summing squares proves the claim. In particular, a
point on this chord within distance $z$ of a binary word is within
distance $z$ of one of the chord endpoints.

Consider first segments varying the same block. If their distance is
less than $d_*$, all fixed encoding blocks and repeated binary levels
agree. For block $2$, the tail agrees and determines the outgoing
edge. For block $1$, the head agrees and determines the incoming edge.
For block $3$, an upward segment has a fixed unequal pair $(E_u,E_v)$,
whereas a downward segment has an equal pair $(E_v,E_v)$; the two types
cannot have the same fixed pair. Within either type that pair determines
the segment, using uniqueness of the incoming edge for the downward
type. There is only one block-$4$ segment.

For two segments varying different blocks, choose points at distance
$z$. Each varying block of a horizontal segment is within distance $z$
of a fixed binary block of the other. By~\eqref{geo:binary-chord}, its
point is within distance $z$ of an endpoint of its original segment.
For the special vertical, its distance from a horizontal segment forces
its height to be at most $z$, so it is within $z$ of its bottom. The
two selected endpoints are therefore within $3z$. Distinct original
endpoints are separated by at least $d_*$: one of their encoding blocks,
binary third levels, or vertical levels differs. Taking $3z<d_*$ forces
the endpoints to coincide.

An unequal ordered pair $(E_u,E_v)$ occurs at the two internal toggle
corners of its own edge only, with the third level distinguishing them.
An endpoint $(E_v,E_v,1,0)$ belongs to the last internal corner of the
unique incoming edge. An endpoint $(E_v,E_v,0,0)$ belongs to at most
the incoming final segment and the outgoing first segment. If both
exist, they are consecutive. At $w_*$ there is no incoming graph edge,
and the special bottom takes exactly that role. The top belongs only
to the special vertical. Thus equality of endpoints of distinct
segments implies consecutiveness. For example $d_{\rm sep}=d_*/4$
is sufficient for the asserted strict-distance argument. The adjacency
components are paths and cycles; any cycle contains at least eight
segments because the graph has no self-edge. Adding a segment at a
path beginning creates no triangle.
\end{proof}

Choose a positive rational $a_0$ such that
\begin{equation}\label{geo:trim-choice}
 2a_0<d_{\rm sep}/2,\qquad4a_0<d_*,\qquad a_0<1/10.
\end{equation}
Trim arclength $a_0$ from both ends of every horizontal segment, whether
or not its outer graph junction is present. At a shared corner $C$ with
incoming and outgoing unit directions $u_0,v_0$, join the trim points by
\begin{align}
 Z(s)&=C-a_0u_0+a_0v_0
             +a_0(\sin s\,u_0-\cos s\,v_0),
             &0\le s\le\pi/2,\label{geo:fillet}\\
 T(s)&=\cos s\,u_0+\sin s\,v_0.\label{geo:fillet-tangent}
\end{align}
The arc has speed $a_0$ in the parameter $s$, and $T$ is its directed
unit tangent. Figure~\ref{fig:geo-fillet} shows this replacement in the
plane of its two orthogonal directions. Use the same fillet at the special
bottom, trimming the vertical there only. Its top is not trimmed. Unmatched horizontal trim
points remain endpoints. Let $\mathcal L$ be the resulting compact union
and set
\begin{equation}\label{geo:endpoint-set}
 \mathcal E=\{\text{unmatched endpoints of }\mathcal L\}
                \setminus\{(H_*,2)\}.
\end{equation}

The first three blocks stay in $[0,1]$: on a fillet the separate varied
block coordinates stay in the intervals traversed by their original
segments. The fourth block is zero except on the vertical and its bottom
fillet; on that fillet it is a repeated scalar in $[0,a_0]$.
Consequently $\langle T,\mathbf e\rangle\le0$ everywhere. Fix an
integer $C_t\ge2/d_*$. For arclength $\lambda$,
\begin{equation}\label{geo:tangent-bounds}
 \norm T_\infty\le C_t,\qquad
 \norm{dT/d\lambda}\le1/a_0,\qquad
 \norm{dT/d\lambda}_\infty\le2C_t/a_0.
\end{equation}
These derivative bounds hold on piece interiors and give the same
Lipschitz bounds across joins because the tangent is continuous there.

\begin{figure}[htbp]
\centering
\setlength{\unitlength}{1pt}
\begin{picture}(290,160)
\thinlines
\multiput(180,30)(-5,0){10}{\line(-1,0){2}}
\multiput(180,30)(0,5){10}{\line(0,1){2}}
\put(180,30){\circle*{3}}
\put(187,24){\makebox(0,0)[l]{$C$}}
\put(130,80){\circle*{3}}
\put(122,87){\makebox(0,0)[r]{$O$}}
\thicklines
\put(20,30){\line(1,0){110}}
\put(62,30){\vector(1,0){25}}
\put(74,18){\makebox(0,0)[t]{$u_0$}}
\qbezier(130,30)(139.946,30)(149.134,33.806)
\qbezier(149.134,33.806)(158.323,37.612)(165.355,44.645)
\qbezier(165.355,44.645)(172.388,51.677)(176.194,60.866)
\qbezier(176.194,60.866)(180,70.054)(180,80)
\put(180,80){\line(0,1){65}}
\put(180,103){\vector(0,1){24}}
\put(170,116){\makebox(0,0)[r]{$v_0$}}
\put(130,30){\circle*{3}}
\put(180,80){\circle*{3}}
\put(130,9){\makebox(0,0)[t]{$C-a_0u_0$}}
\put(188,80){\makebox(0,0)[l]{$C+a_0v_0$}}
\end{picture}
\caption{Rounding one orthogonal corner. The dashed legs are removed;
the directed path follows the retained straight portions and the fillet
from $C-a_0u_0$ to $C+a_0v_0$. Its center is
$O=C-a_0u_0+a_0v_0$ and its radius is $a_0$. This local geometry lies
in the orthonormal $(u_0,v_0)$ plane.}
\label{fig:geo-fillet}
\end{figure}

The numerical procedure will choose a curve point by minimizing estimated
distances. We therefore need more than a unique projection: a curve point
with small squared-distance excess must be close to the true projection,
and its tangent close to the projection's tangent, in every coordinate.
The next lemma gives these bounds before any candidate is selected.

\begin{lemma}[Local chains and strong excess]\label{geo:projection}
There is $\rho_{\rm loc}>0$, depending only on $d_*,a_0$, such that
two points $Z,Z'\in\mathcal L$ with $\norm{Z-Z'}<\rho_{\rm loc}$
are joined along one local straight or single-corner chain. Its length
$\Delta$ satisfies $\Delta\le\sqrt2\norm{Z-Z'}$.
For every sufficiently small $\tau>0$, if
\begin{equation}\label{geo:tube-hypothesis}
 r=\dist(X,\mathcal L)\le4\tau,\qquad
 \dist(X,\mathcal E\cup\{(H_*,2)\})>10\tau,
\end{equation}
there is a unique closest point $Z$, which is interior and satisfies
$\langle X-Z,T(Z)\rangle=0$. If $Z'\in\mathcal L$ has squared excess
$g=\norm{X-Z'}^2-r^2\le\tau^2$ and signed arclength difference $s'$
from $Z$ on that chain, then
\begin{equation}\label{geo:strong-excess}
 g\ge (s')^2/2.
\end{equation}
In particular,
\begin{align}
 \norm{Z'-Z}&\le\sqrt{2g},&
 \norm{Z'-Z}_\infty&\le C_t\sqrt{2g},\label{geo:point-error}\\
 \norm{T(Z')-T(Z)}&\le\sqrt{2g}/a_0,&
 \norm{T(Z')-T(Z)}_\infty&\le2C_t\sqrt{2g}/a_0.
 \label{geo:tangent-error}
\end{align}
The tangent at an endpoint competitor is its one-sided tangent.
\end{lemma}
\begin{proof}
Support a straight point by its original segment, and an arc point by
the two segments of its corner. An arc point is within $a_0$ of each
supporting segment. Choose $\rho_{\rm loc}+2a_0<d_{\rm sep}$.
Every across-point pair of supports is then identical or consecutive.
Each point's support set is itself a clique, so their union is a clique
and contains at most two segments by Lemma~\ref{geo:separation}. In the
one-segment case both points lie on its retained straight. In the
two-segment case they lie on the two retained straights and their
shared fillet. Opposite-end trims are disjoint by~\eqref{geo:trim-choice}.
This also rules out unintended identifications or intersections.

Along a two-direction chain, all tangent coefficients in its orthogonal
directions $u_0,v_0$ are nonnegative, and
$\langle T,u_0+v_0\rangle\ge1$. Integrating in the directed order gives
$\Delta\le\langle Z'-Z,u_0+v_0\rangle\le\sqrt2\norm{Z'-Z}$.
Reverse the order if necessary. A straight chain has equality between
length and endpoint distance.

Compactness gives a minimizer. The endpoint assumption makes every
minimizer interior, so differentiating squared distance along its
$C^1$ path gives the asserted perpendicularity. A competitor of the
stated excess satisfies
\[
 \norm{Z'-Z}\le4\tau+\sqrt{17}\tau<9\tau.
\]
Require $9\tau<\rho_{\rm loc}$. The local chain then has
$|s'|<13\tau$, and integration of~\eqref{geo:tangent-bounds} gives
\[
 Z'-Z=s'T(Z)+R,\qquad \norm R\le(s')^2/(2a_0).
\]
For $|s'|/(2a_0)<1$, expansion of squared distance and perpendicularity
therefore imply
\begin{align*}
 g&=\norm{Z'-Z}^2-2\langle X-Z,Z'-Z\rangle\\
  &\ge(s')^2\left[
      \left(1-\frac{|s'|}{2a_0}\right)^2-\frac r{a_0}\right]
   \ge(s')^2/2,
\end{align*}
where, for example, $\tau/a_0\le1/64$ makes the last inequality hold.
Applying this to a second minimizer proves uniqueness. Integrating the
unit tangent, its infinity bound, and its two Lipschitz bounds over
length $|s'|\le\sqrt{2g}$ proves~\eqref{geo:point-error} and
\eqref{geo:tangent-error}. All estimates are independent of $m$.
\end{proof}

At an unmatched beginning, the trimmed segment varies only block $2$;
at an unmatched ending it varies only block $3$. Their first blocks
are therefore exactly the encoding of the missing-neighbor vertex.
There are no isolated vertices requiring a separate convention. The
source's horizontal beginning is joined to the special vertical and
is not an endpoint of $\mathcal L$. This proves the endpoint assertion
in Theorem~\ref{thm:geometry}. If $P\in\mathcal E$ and $\norm{X-P}\le10\tau$, every wrongly
rounded first coordinate costs at least $1/4$ in squared error, proving
the Hamming bound in that theorem.

\begin{remark}[Local Lipschitz continuity of projection]
\label{geo:projection-lipschitz}
Suppose both $X$ and $Y$ satisfy~\eqref{geo:tube-hypothesis}, and
$\norm{X-Y}\le\tau/20$. Write $P,Q$ for their respective projections
onto $\mathcal L$ and $r_X,r_Y$ for their distances to $\mathcal L$.
With the same sufficiently small choice of $\tau$, one has
\begin{equation}\label{geo:projection-lipschitz-bound}
 \left(1-\frac{r_X+r_Y}{a_0}\right)\norm{P-Q}^2
 \le \langle X-Y,P-Q\rangle
 \le \norm{X-Y}\norm{P-Q}.
\end{equation}
In particular, the projection is locally $2$-Lipschitz on this tube.
This is a local assertion away from the specified endpoints.

Indeed, the conclusion is immediate if $P=Q$. Otherwise,
$\norm{P-Q}\le r_X+\norm{X-Y}+r_Y\le(8+1/20)\tau<\rho_{\rm loc}$.
Lemma~\ref{geo:projection} therefore puts $P,Q$ on one local chain,
whose arclength $s$ satisfies $s^2\le2\norm{P-Q}^2$.
Integrating~\eqref{geo:tangent-bounds} from either end expresses its
chord as the signed length times the tangent there, with remainder
of norm at most $s^2/(2a_0)$. Perpendicularity at $P$ and $Q$ gives
\[
 \big|\langle X-P,P-Q\rangle\big|
       \le\frac{r_Xs^2}{2a_0},\qquad
 \big|\langle Y-Q,P-Q\rangle\big|
       \le\frac{r_Ys^2}{2a_0}.
\]
Expanding $X-Y=(P-Q)+(X-P)-(Y-Q)$ proves the first inequality
in~\eqref{geo:projection-lipschitz-bound}; Cauchy--Schwarz proves
the second. Finally $r_X+r_Y\le8\tau$ and $\tau/a_0\le1/64$
make its left coefficient at least $7/8$, and in particular at least
$1/2$, proving the stated factor $2$.
\end{remark}

\subsection{A field with uniform displacement after clipping}

For the inside prescription, away from the desired endpoint neighborhood,
successive radial intervals interpolate from the directed tangent through
the inward normal and reverse tangent to the fixed fourth-block direction.
The four directional stages in~\eqref{geo:Phi} parallel
\cite[Equation~(1)]{Rubinstein2016Settling}; the projection and clipping estimates
below establish the bounds needed for the present curve.
A separate prescription above fourth-block mean level $1/2$ handles the
distinguished vertical. We check that the prescriptions agree and that
clipping preserves positive displacement away from the desired endpoints.
The normalized tube bound permits individual coordinates of $X$ to be
far from the curve, so raw field length alone does not prove this last
claim. We will bound the fraction of coordinates where clipping can
remove motion, and then handle the upper face separately.
\label{geo:field}

Fix $\tau$ satisfying Lemma~\ref{geo:projection}, and impose also
\begin{equation}\label{geo:tau-choice}
 \tau<1/100,\qquad6(C_t+1)\tau\le1/4.
\end{equation}
All further fixed validation or recovery upper bounds on $\tau$ are
imposed now. For $X=(X_1,X_2,X_3,X_4)$ put
$\mu=m^{-1}\sum_j(X_4)_j$. For a point $Z$ and a unit tangent $T$ let
$r=\norm{X-Z}$ and define
\begin{equation}\label{geo:Phi}
 \Phi(X;Z,T)=
 \begin{cases}
 (1-r/\tau)T+(Z-X)/\tau,&0\le r\le\tau,\\[2pt]
 -(r/\tau-1)T+\dfrac{2\tau-r}{\tau r}(Z-X),
                       &\tau\le r\le2\tau,\\[5pt]
 -\dfrac{3\tau-r}{\tau}T+\dfrac{r-2\tau}{\tau}\mathbf e,
                       &2\tau\le r\le3\tau,\\[5pt]
 \mathbf e,&r\ge3\tau.
 \end{cases}
\end{equation}
The values at shared boundaries agree. There is no reciprocal radius
in the formula used near zero.

For $\mu\le1/2$ and $\dist(X,\mathcal E)>10\tau$, use
$v(X)=\Phi(X;Z,T(Z))$ when $\dist(X,\mathcal L)\le3\tau$, and
$v(X)=\mathbf e$ otherwise. The top is at distance at least $3/2$
from this region, by the difference of fourth-block means; hence the
projection lemma applies. At the remaining inside points put $v(X)=0$.
For $\mu\ge1/2$, the default upward motion can be lost at the upper
face. We add motion toward the distinguished vertical, with coefficient
zero at the mean interface. Close to that vertical the tangent direction
is downward, including at its top. Set
\begin{equation}\label{geo:outside}
 \begin{aligned}
 Z_\mu&=(H_*,\mu),\qquad r_\mu=\norm{X-Z_\mu},\qquad
 h_\mu=\min\{1,2(\mu-1/2)\},\\
 v(X)&=\Phi(X;Z_\mu,-\mathbf e)+h_\mu(Z_\mu-X).
 \end{aligned}
\end{equation}

\begin{lemma}[Agreement at the mean interface]\label{geo:interface}
In $2/5<\mu<3/5$, the inside prescription extended to this slab equals
$\Phi(X;Z_\mu,-\mathbf e)$. The outside extension, using
$h_\mu=\max\{0,\min\{1,2(\mu-1/2)\}\}$, differs from it in norm by
at most $12|\mu-1/2|$ and in each component by at most $6|\mu-1/2|$.
Thus the definitions agree at $\mu=1/2$.
\end{lemma}
\begin{proof}
Every nonvertical piece, including the bottom fillet, has fourth height
at most $a_0<1/10$. Its distance to a point in the slab exceeds $3/10$,
by the mean inequality $\norm{X_4-z\mathbf1}\ge|\mu-z|$.
It cannot be in the $4\tau$ tube. The retained vertical contains all
heights of this slab, and
\[
 m^{-1}\sum_j((X_4)_j-t)^2
 =m^{-1}\sum_j((X_4)_j-\mu)^2+(t-\mu)^2
\]
shows that its nearest line point is $Z_\mu$. For $T=-\mathbf e$,
the last annulus of~\eqref{geo:Phi} is already identically
$\mathbf e$ at every $r\ge2\tau$. Consequently the same equality
holds even where the closest curve piece is far away and the inside
field is default. The desired endpoints have fourth height zero and
the top has height two, so neither endpoint exception occurs in the
slab. Each coordinate of $Z_\mu-X$ has magnitude at most three and
its norm is at most six. Multiplication by
$h_\mu\le2|\mu-1/2|$ proves both bounds.
\end{proof}

The projection is continuous wherever it is unique in the tube:
if $X_n\to X$, compactness makes every subsequential limit of their
projections a minimizer for $X$, hence the unique minimizer. Its tangent
is continuous along the local chains. Formula~\eqref{geo:Phi} matches
the default at $3\tau$, and Lemma~\ref{geo:interface} matches the mean
interface. This proves precisely the continuity claimed in
Theorem~\ref{thm:geometry}.

\begin{proof}[Proof of the displacement assertion in Theorem~\ref{thm:geometry}]
Choose $\beta_0>0$ satisfying all the upper bounds on $\beta$ used below.
\paragraph{Inside the tube and in the inside default region.}
First suppose $\mu\le1/2$ and $r=\dist(X,\mathcal L)\le3\tau$.
For $r\le\tau$ the raw field is the convex combination of $T$ and
the inward unit normal with coefficients $1-r/\tau,r/\tau$.
For $\tau\le r\le2\tau$ it is the convex combination of $-T$ and
the inward unit normal with coefficients $r/\tau-1,2-r/\tau$.
The two directions in each combination are orthogonal. At $r=0$
the field is $T$. In the third annulus the two unit directions are
$-T,\mathbf e$, whose inner product is nonnegative. Therefore in all
three annuli $\norm{v(X)}\ge1/\sqrt2$.

Write $v=b+w(Z-X)$, where $\norm b_\infty\le C_t+1$ and
$0\le w\le1/\tau$. All coordinates of $Z$ lie in $[0,1]$: for a
vertical projection its height is at most $\mu+3\tau<1$, and all
other pieces already have this property. Require
$\beta/\tau\le1/2$ and $\beta(C_t+1)\le1/4$.
The preliminary point $Y=X+\beta w(Z-X)$ lies in the box. If
$|X_i-Z_i|<1/2$, then $-1/2<Y_i<3/2$, and the additional displacement
$\beta b_i$ causes no clipping. The number of remaining coordinates,
divided by $m$, is at most $4r^2\le36\tau^2$. At each, the clipping
loss is at most $\beta(C_t+1)$ because $Y$ was in the box. The norm of
all loss is at most $6\tau\beta(C_t+1)\le\beta/4$. Subtracting it
from the raw displacement leaves at least $\beta/4$.

In the inside default case, let $p$ be the fraction of fourth coordinates
at most $7/4$. Since all entries are at least $-1$,
\[
 1/2\ge\mu\ge-p+(7/4)(1-p),\qquad p\ge5/11.
\]
For $\beta\le1/4$ these coordinates retain the full upward displacement,
giving norm at least $\beta\sqrt{5/11}$.

\paragraph{Above the mean interface.}
Now suppose $\mu\ge1/2$, and abbreviate $Z=Z_\mu$, $r=r_\mu$.
The fourth coordinates of $Z-X$ have mean zero, so they are orthogonal
to $\mathbf e$. Write $v=b\mathbf e+w(Z-X)$, with the scalar table
\begin{equation}\label{geo:outside-coefficients}
\begin{array}{c|cc}
\text{radius}&b&w\\\hline
0\le r\le\tau&-1+r/\tau&1/\tau+h_\mu\\
\tau\le r\le2\tau&r/\tau-1&2/r-1/\tau+h_\mu\\
r\ge2\tau&1&h_\mu
\end{array}
\end{equation}
and hence
\begin{equation}\label{geo:outside-norm}
 |b|\le1,\quad0\le w\le w_{\max}:=1/\tau+1,
 \quad b^2+w^2r^2\ge1/2.
\end{equation}
The last assertion follows from the two orthogonal components of the
first two radial regimes; adding $h_\mu$ increases the nonnegative
normal coefficient. In the last regime $b=1$. Require
$\beta w_{\max}\le1/2$. The move toward $Z$ is a convex move inside
the box, and horizontal coordinates undergo no further move or clipping.

If $b<0$, then $r<\tau$. The subsequent downward step cannot clip at
the upper boundary. Where $|(X_4)_j-\mu|<1/2$, the preliminary fourth
coordinate is nonnegative since $\mu\ge1/2$, and the following step
is at least $-\beta\ge-1$. Thus lower clipping can occur only on a
fraction at most $4\tau^2$ of the fourth coordinates. Its loss on each
is at most $\beta$, giving norm loss at most $2\tau\beta$. By
\eqref{geo:outside-norm} this leaves at least $\beta/4$.

It remains to treat $b\ge0$. Set
\[
 \sigma^2=m^{-1}\sum_j((X_4)_j-\mu)^2,\qquad
 H=\norm{(X-Z)_{1,2,3}},\qquad r^2=H^2+\sigma^2,
\]
and choose a positive rational
\begin{equation}\label{geo:eta-choice}
 \eta\le\min\left\{\frac\tau2,\frac1{4w_{\max}},
           \frac12\min\left(\frac12,\frac1{4w_{\max}}\right)\right\}.
\end{equation}
If $\sigma\ge\eta$, a uniform fourth coordinate satisfies
\[
 \sigma^2\le3\E|(X_4)_j-\mu|
        =6\E(\mu-(X_4)_j)_+.
\]
Put $c=\eta^2/12$. Since the positive part is at most three,
$\E(\mu-(X_4)_j)_+\ge\eta^2/6$ implies that a fraction at least
$\eta^2/36$ have $\mu-(X_4)_j\ge c$: otherwise its expectation
would be less than $c+3\eta^2/36=\eta^2/6$.
On those coordinates the room to the upper boundary is at least $c$.
Also $b+wc\ge c/12$. Indeed~\eqref{geo:outside-norm} and $r\le6$
imply that either $b\ge1/2$ or $w\ge1/12$; $c<1$ handles the first
alternative. For $\beta\le1$, clipping therefore retains displacement
at least $\beta c/12$ on that fraction. Its norm is at least
$\beta\eta^3/864$.

Suppose instead $\sigma<\eta$. If $wH\ge\tau/4$, the horizontal
displacement, which is not clipped, is at least $\beta\tau/4$.
Otherwise~\eqref{geo:outside-norm} gives
\[
 b^2\ge\frac12-w^2H^2-w^2\sigma^2
       \ge\frac12-\frac{\tau^2}{16}-w_{\max}^2\eta^2
       \ge\frac38,
\]
so $b\ge1/2$. If $\mu\le1$, let
$d_0=\min\{1/2,1/(4w_{\max})\}$. By~\eqref{geo:eta-choice}, at
least three quarters of the fourth coordinates have deviation from
$\mu$ at most $d_0$. Their upper room is at least $1/2$, and their
velocity is at least $b-w_{\max}d_0\ge1/4$. They retain at least
$\beta/4$ for $\beta\le1$, giving norm at least
$\sqrt3\beta/8>\beta/8$. If $\mu>1$, then $h_\mu=1$, $w\ge1$,
and $b\ge0$ in~\eqref{geo:outside-coefficients} requires $r\ge\tau$.
It follows that
\[
 wH\ge H\ge\sqrt{\tau^2-\eta^2}\ge\sqrt3\tau/2,
\]
contrary to the remaining case assumption. This exhausts all outside
points, including the upper face. At $(H_*,2)$, $r=0$ and the formula
gives $v=-\mathbf e$ exactly.

Take positive rational
$c_*\le\min\{1/8,\tau/4,\eta^3/864\}$ and positive rational
$\beta_0$ satisfying all restrictions just made. All coordinate
differences used in the formulas have magnitude at most three, so any
integer $B_1\ge C_t+6+3/\tau$ satisfies~\eqref{geo:bounded-field},
including the zero convention near inside endpoints. These choices
depend only on the previously fixed geometric constants. This completes
the proof of Theorem~\ref{thm:geometry}.
\end{proof}

\subsection{The local-access contract}
\label{geo:access}

We next state the information needed to approximate this particular
field. All probability statements below are for a fixed physical vector
$X$, averaged over a uniform physical component $I\in[4m]$ and the
public trial randomness. Write $J=1+((I-1)\bmod m)$ for its word
coordinate. Fix a common threshold $\xi\in[2/5,3/5]$ and round the
first two blocks at $\xi$. The two resulting binary words $y_1,y_2$
are fixed, although the procedures only read requested coordinates.
Any scalar-quantization offset may also be fixed throughout the following
contract and accuracy analysis.

\begin{definition}[Admissible local access]\label{geo:access-contract}
An \emph{oracle access implementation} satisfies parts (a)--(c) below
and the following operational condition. At its declared oracle level,
each supplied bit or decision routine uses a bounded tree of completed
Boolean answers, apart from public index computations. Every input,
including an input outside a correctness promise, gives Boolean answers
or public defaults. After its requested accuracies are fixed, the
oracle tree and its symbolic query roles are bounded independently of
$m$ and remain fixed under refinement of threshold-offset grids.
These conditions impose no constant bound on the physical circuit
realizing a completed oracle answer. Part (d) is an additional
\emph{physical-read hypothesis}, required only when a conclusion about
the fully expanded physical reads is explicitly invoked. All constants
are uniform in $m$, $X$, and admissible fixed offsets. Geometric and
validation radii and the reference rule below are fixed before
reliability and numerical accuracy choices. Operational, sample,
repetition, and read-density bounds may depend on the requested fixed
accuracies as well as the encoding construction.
\begin{enumerate}[label=\textup{(\alph*)},leftmargin=*]
\item\label{geo:access-validation}
Each seed has a validation decision whose failure probability can be
made any prescribed fixed positive number. For fixed $y_s$ there is
either no valid reference or a single fixed vertex $u_s$ such that, off
the validation failure event, acceptance identifies $u_s$. This
reference is independent of the actual random acceptance decision.
It is also independent of requested confidence and numerical accuracy.
If $\norm{X_s-E_u}\le5\tau$, off that failure event the seed is
accepted with reference $u$. Intermediate seeds with a proper reference
may be accepted or rejected. An accepted seed is used only through
local routines, not through a globally decoded vertex name.
\item\label{geo:access-bits}
For any accepted reference, incident-edge flags can be supplied with
arbitrarily small fixed error, and bits of its own and its actual
immediate neighbors' encodings can be supplied with any prescribed
fixed average error. The same bit guarantee holds for $E_{w_*}$ and
the encoding of its actual successor, even if neither seed is accepted.
These special words may be public or virtual.

The bit guarantee is unconditioned for each fixed relevant symbolic
role and branch: a request at a whole-word coordinate with density at
most a fixed constant times $1/m$ has the asserted error under its
given request law. It is not a guarantee conditional on selecting that
branch. Requests to different word operands may use the same complete
coordinate. The guarantee is required for the matching coordinate $J$
and for all preallocated statistical coordinates below. A finite union
of role or fixed block-type alternatives is allowed.
\item\label{geo:access-statistics}
For every finite family of bounded functions $f:[m]\to\mathbb R$
fixed by $X$, the rounded seeds' possible actual references, and their
actual candidate data, one can choose fixed sample counts so that each
sample average estimates $m^{-1}\sum_jf(j)$ to any prescribed fixed
additive tolerance with any prescribed fixed failure probability.
The law is again unconditioned under $(I,\text{trial})$, for each fixed
relevant role alternative. Every individual sampled word coordinate has
uniform marginal, or the prescribed bounded-density marginal supported
by the same bit-access guarantee. Corresponding word operands and
physical blocks are evaluated at matching coordinates.

Independent uniform samples satisfy this requirement. Other sample
schemes are permitted if their concentration statement holds for the
fixed underlying functions in this sense; samples and different batches
need not be independent after averaging over $I$.
\item\label{geo:access-finite}
\emph{Additional physical-read hypothesis.} After a fixed accuracy has
been specified, every supplied routine has a realization with a bounded
number of primitive physical read answers and bounded Boolean
processing, apart from public index computations. The bound is
independent of $m$. On every fixed branch of its complete physical
answer tree, including branches with wrong earlier answers, the
unconditioned marginal of each raw address in a physical domain of
size $M$ is at most $C/M$, with a fixed $C$. A branch that uses public
data need not make a physical read. These bounds hold
uniformly in admissible offsets, whose refinement changes comparison
constants and not address primitives or numbers of hidden-answer slots.
On each fixed symbolic answer branch, the address rule itself is
independent of the numerical values of these offsets.
\end{enumerate}
\end{definition}

Part~\ref{geo:access-bits} controls accuracy of a completed logical call
on a fixed valid reference. The optional
part~\ref{geo:access-finite} additionally controls physical reads on
every wired branch, whether or not it is accurate or executed. The
inner implementation of Section~\ref{sec:inner-gadget} satisfies this
additional hypothesis. The outer construction invokes only the oracle
contract: its ideal functional answers are subsequently replaced by
gadgets of polynomial size in the port length. Its fixed direct-read
count $H_{\rm out}$ and functional-invocation count $B_{\rm out}$,
followed by Theorem~\ref{thm:port-gadget} and
Lemma~\ref{out:routing}, account separately for that physical
substitution. We do not assert a size-independent physical read count
for the expanded outer construction. Neither application requires
per-coordinate correctness.

The physical completeness condition in part~\ref{geo:access-validation}
is obtained from a binary completeness radius $\rho_{\rm comp}>0$ by
the initial choice
\begin{equation}\label{geo:rounding-radius}
 \frac{(5\tau)^2}{(2/5)^2}<\rho_{\rm comp}.
\end{equation}
Indeed a rounded error relative to a binary coordinate costs squared
physical error at least $(2/5)^2$ for every allowed $\xi$. An
acceptance radius smaller than half the encoding's relative Hamming
distance makes its reference unique. In the inner application this
reference must lie on the path for the fixed true input tuple, exactly
as established by the validation in Section~\ref{sec:inner-gadget}.

\begin{remark}[Whole-word coordinates in the inner application]
A matching target is a whole encoding coordinate. In a product block
called $Z$, that coordinate includes both internal labels of the product.
Using the same whole coordinate in $E_u$ and $E_v$ preserves the supplied
joint request law; it does not replace the two internal labels of either
word by a diagonal pair.

For mixed public and variable products in the virtual source, the joint
density bound gives the marginal bound
\[
 p(x,y)\le\frac{C}{|A||B|}
 \quad\Longrightarrow\quad
 \sum_{y\in B}p(x,y)\le\frac C{|A|}.
\]
The product becomes an additive form on the variable operand. That form
may depend on the public operand and the discarded coordinate, but it
does not use that coordinate to select a new variable address.
Same-positive-tag products retain the whole pair, as required by
Lemma~\ref{lem:inner-binary-interface}. Lemma~\ref{ig:virtual-source}
proves this transport from the actual source and first successor to the
fixed dummy ports. Thus their virtual supply has exactly the
whole-coordinate access required by the contract.
\end{remark}

\begin{lemma}[Geometric procedures]\label{lem:geometric-procedures}
Fix the geometric constants of Theorem~\ref{thm:geometry} and an
oracle access implementation satisfying parts (a)--(c) and the
operational condition of Definition~\ref{geo:access-contract}.
For every fixed $a>0$ there is a finite local numerical schema with
properties (i)--(iii) below. Property (iv) additionally assumes the
physical-read hypothesis~\ref{geo:access-finite}.
\begin{enumerate}[label=\textup{(\roman*)},leftmargin=*]
\item It returns a signed value $\widehat v_I(X)$ on a fixed dyadic grid
inside $[-B_1,B_1]$. It terminates and is bounded on every input and
every trial. For every fixed $X$ with $\dist(X,\mathcal E)>10\tau$,
and conditional on each fixed admissible binary threshold and
scalar-quantization offset,
\begin{equation}\label{geo:average-accuracy}
 4\E_{I,\,\mathrm{trial}}
       |\widehat v_I(X)-v_I(X)|^2\le a^2.
\end{equation}
\item At fixed rounded seeds, correct decisions refer to at most eight
true-word roles and 32 actual piece roles. Meshes, scalar tolerances,
sample counts, and all statistical and final matching bit requests
are fixed before choosing the smaller bit-error budgets, and all final
matching requests are allocated before numerical minimization.
\item Apart from the supplied routines and public index computations,
the schema uses finitely many scalar comparisons and bounded Boolean
lookup. Its only word-coordinate address primitives are the matching
coordinate and the preallocated statistical coordinates. This remains
true on every hypothetical wrong or default branch.
If the oracle-level public spaces, probabilities, routine descriptions,
and address computations are uniformly enumerable with polynomial
support, time, and rational-probability bit complexity in a caller's
size parameter, the composed oracle-level schema has those properties
as well. This does not include physical expansion of an oracle call.
\item If in addition the physical-read
hypothesis~\ref{geo:access-finite} holds, then, once all accuracies and
implementation counts have been fixed, the fully unfolded physical
primitive tree has a fixed number $H$ of possible read slots. For each
physical domain of size $M$, its unconditioned
expected occurrence count at any fixed position, averaged over $I$
and public trials, is at most $C/M$, for a fixed sum $C$ of slot
density bounds. Both $H,C$ remain fixed under subsequent refinement
of threshold-offset grids. Under the corresponding polynomial
enumerability premise for these physical realizations, the fully
expanded physical schema also has the enumerability properties in
part (iii).
\end{enumerate}
Every choice concerns the unscaled field. In particular none depends
on a subsequently chosen $0<\beta\le\beta_0$.
\end{lemma}

\subsection{A finite list covering every close piece}
\label{geo:candidates}

From a seed reference $u$, reserve the encoding roles for $u$, its
predecessor, and its successor. Reserve its two possible incident edge
roles. For each actual incident edge list the four retained straight
segments and its three internal fillets. The trim convention uses the
same $a_0$ at both outer ends regardless of whether those ends are
shared, so these seven pieces need only their two endpoint words.
When both graph edges at $u$ exist, also list the graph-junction fillet
at $u$. The incoming direction there is the fixed downward third-block
direction, so the junction needs only $E_u$, the outgoing neighbor,
and the two flags. Apply this procedure to both seeds.

Always add the special retained vertical and its bottom fillet using
the distinguished source and its actual successor. These two pieces
are present even on a trial with no accepted seed. The ordinary
junction rule requires both graph edges and therefore never invents
another junction at the source, which has no graph predecessor.
There is no need to recognize the source by comparing hidden names.

This is the following explicit reservation schedule; duplicate roles
do no harm.
\begin{center}
\begin{tabular}{@{}lrr@{}}
\toprule
Reserved roles&Word roles&Piece roles\\
\midrule
Self, predecessor, successor for two seeds&6&0\\
Four incident edges, seven pieces each&0&28\\
Two ordinary junctions&0&2\\
Distinguished source and actual successor&2&0\\
Special vertical and bottom fillet&0&2\\
\midrule
Total upper bound&8&32\\
\bottomrule
\end{tabular}
\end{center}
On good validation and flag answers, the admitted roles are actual
pieces. A gray acceptance merely chooses whether to include one fixed
reference's list. For fixed $X,\xi$, all possible underlying actual
words and pieces in the table are fixed independently of later accuracy
choices and numerical decisions.

\begin{lemma}[Two-seed coverage]\label{geo:coverage}
Every non-special piece has either its first or its second block fixed
to an encoding of a vertex whose seed generates it. Consequently,
on the strong-completeness and flag good event, every piece at distance
at most $4\tau$ from $X$ is included in the reserved candidate list.
\end{lemma}
\begin{proof}
For $u\to v$, its first segment, first fillet, and upward toggle have
first block $E_u$. The fillet from that toggle to the block-$1$
segment, the block-$1$ segment, and its fillet to the downward toggle
have second block $E_v$. The downward toggle has both blocks $E_v$.
An ordinary graph-junction fillet at $u$ changes only blocks $3$ and
$2$, so its first block is $E_u$. In each case the stated vertex
generates the piece from its incident data. If the piece is within
$4\tau$ of $X$, its fixed block is within $4\tau$ of the corresponding
physical block. Part~\ref{geo:access-validation} of the contract
therefore accepts that seed with its true reference. Correct flags
include the piece. The two special pieces are included unconditionally.
Gray acceptance cannot remove a piece required by this argument; it
can only add or remove other actual pieces.
\end{proof}

\subsection{Numerical data and the order of their estimation}
\label{geo:numerical}

Put a finite parameter mesh, including endpoints, on each straight
piece and each angle interval $[0,\pi/2]$ of a fillet. All original
lengths are at most two, so a sufficiently fine fixed mesh gives
maximum arclength spacing at most any prescribed fixed $D>0$.
For the proof, a mesh entry means its exact point $Z'$ and its exact
tangent on its underlying actual piece. The implementation computes
approximations to these entries.

At most five endpoint-word distance statistics are needed: the four
ordinary edge roles and the special edge. Each estimates
\begin{equation}\label{geo:length-statistic}
 \ell_{uv}^2=m^{-1}\sum_{j=1}^m(E_u(j)-E_v(j))^2.
\end{equation}
Use the same coordinate in the two words. The true nonconstant
lengths are in $[d_*,1]$; clamp their estimated square roots to
$[d_*/2,2]$. The formulas for trim points, straight tangents, and
\eqref{geo:fillet}--\eqref{geo:fillet-tangent} then have only bounded
denominators separated from zero. Their coordinates are uniformly
continuous functions of the endpoint bits at that coordinate, the
lengths, and the finitely many fixed trigonometric mesh constants.
Thus correct endpoint bits and sufficiently accurate length and public
constant approximations give any desired component error, uniformly
in the coordinate and in $m$. A component error at most $z$ everywhere
gives full-vector norm error at most $2z$.

Also estimate the fourth-block mean $\mu$. For each of the at most
32 meshed roles, estimate each exact mesh point's squared distance by
the sum of its four block means,
\begin{equation}\label{geo:mesh-statistic}
 \norm{X-Z'}^2
   =\sum_{b=1}^4m^{-1}\sum_j(X_b(j)-Z'_b(j))^2.
\end{equation}
Finally reserve the corresponding outside-radius statistic for the
exact point $Z_\mu$. Physical samples and endpoint-bit samples in a
summand use the same $j$. One may use separate batches for every
statistic and alternative role. Before a numerical winner is selected,
reserve the final matching-coordinate requests for all eight word
roles, as well as the matching physical coordinate. Inactive requests
may later be omitted or use public defaults.

These data define the numerical estimator. Use the estimated mean to
choose the inside or outside prescription, with a fixed rule at equality.
For the inside prescription, select an admitted mesh entry with minimum
estimated squared distance, again with a fixed rule for ties. Clamp that
estimate at zero and take its square root. Substitute this scalar radius,
the computed tangent component, and the computed component of $Z'-X$
into the component formula~\eqref{geo:Phi}; its last regime gives the
default value. The radius is supplied by the statistic, not by a
full-vector norm computation. For the outside prescription, clamp the
outside squared-radius estimate at zero and take its square root. Use
this radius, the estimated mean, and the special matching bits in
\eqref{geo:outside}. Bounded clamping and final dyadic rounding
will make this a bounded procedure on every trial. We next prove that
the estimates are accurate and that the resulting component is stable
under approximate minimization and side selection.

Here is a precise reason that estimates based on earlier approximate
lengths still satisfy the required concentration statement. Fix $X$, the
rounded seeds, and one possible actual reference assignment. The exact
functions in~\eqref{geo:length-statistic}, $\mu$,
\eqref{geo:mesh-statistic}, and the outside-radius statistic are fixed
functions of the sample index; they do not depend on an estimated
length or a numerical winner. They have bounded ranges: a squared
physical difference is at most nine, and the sum of four such terms
is at most 36. Apply part~\ref{geo:access-statistics} to these exact
functions. On the event of correct sampled bits, earlier good length
estimates and good numerical constants perturb every implemented
summand by a uniformly small amount. For example if $|z-z'|\le t$
and $x,z,z'\in[-2,3]$, then
\[
 |(x-z)^2-(x-z')^2|
   =|z-z'||2x-z-z'|\le10t.
\]
Clamps keep all approximations in one such fixed range. The analogous
bound for replacing a physical read by its quantized representative
is the same. For the outside point, a good mean estimate changes all
fourth coordinates by at most its scalar error; the other coordinates
use special-word bits. It follows that concentration of the exact
statistics yields accuracy of all implemented statistics after these
small deterministic perturbations. No assertion conditioned on the
earlier estimates is being used.

For independent uniform sampling, this concentration needs no external
input. If $|f|\le M$, independence gives
$\operatorname{Var}(q^{-1}\sum_{j=1}^q f(K_j))\le M^2/q$.
Chebyshev's inequality gives failure probability at most
$M^2/(qz^2)$ at tolerance $z$; a union bound handles the fixed finite
list. In the outer application the batches are instead stationary
under the output-coordinate average. The concentration calculation in
Section~\ref{sec:outer-composition} supplies the same required small
failure probabilities for these fixed functions. The geometry uses
that conclusion, and does not condition on a gray acceptance pattern,
a realized mesh winner, or a particular output coordinate.

\begin{lemma}[Uniform component stability]\label{geo:stability}
Fix a positive desired component error $e$. Mesh resolution,
quantization width, scalar tolerances, and precision of public constants
can be chosen so that, whenever all relevant decisions, exact scalar
statistics, and requested bits are good, the local procedure's component
error is at most $e$ at every $X$ outside the endpoint neighborhood.
The bound holds for every output coordinate, even if the approximate
winner depends on that coordinate. It also holds with arbitrary tie
choices between numerically minimizing entries.
\end{lemma}
\begin{proof}
For the inside prescription, consider the selected mesh entry.
Let $z>0$ bound every squared-distance error, including the parameter,
bit, quantization, and scalar-statistic errors just discussed. Suppose
first that $r=\dist(X,\mathcal L)\le3\tau$. By
Lemma~\ref{geo:coverage}, a minimizing piece is present. A mesh point
at arclength at most $D$ from its projection has squared excess at most
$2rD+D^2$. Every admitted exact entry is on the actual curve, so none
has negative excess. The exact point $Z'$ underlying the approximate
winner therefore has
\begin{equation}\label{geo:winner-excess}
 0\le g=\norm{X-Z'}^2-r^2\le2rD+D^2+2z.
\end{equation}
Choose $D,z$ first small enough that the right side is at most
$\tau^2$, and then still smaller as required. Equations
\eqref{geo:point-error}--\eqref{geo:tangent-error} show that both its
point and tangent approach those of the true projection in infinity
norm, uniformly. Correct endpoint bits at the final matching coordinate
also make its computed point and tangent close to their exact values.
Clamping a squared-radius estimate at zero and taking square root is
uniformly continuous on the fixed bounded interval. Consequently its
estimated radius approaches $r$.

Each component of~\eqref{geo:Phi} is a bounded continuous expression
in that radius, the tangent component, and the components of $Z-X$.
The formulas match at $\tau,2\tau,3\tau$, and the only reciprocal
radius is protected by $r\ge\tau$ in its regime. Safe clamping can be
used at approximate regime boundaries. Uniform continuity on these
fixed compact ranges proves the desired component error. Crucially,
the infinity-norm bounds hold for every good candidate satisfying
\eqref{geo:winner-excess}; they therefore also hold when the chosen
candidate is correlated with the requested coordinate.

If the true inside radius is $r>3\tau$, the true field is default.
Every actual candidate entry has distance at least $r$. Good squared
distance estimates can yield a nondefault computed expression only
arbitrarily near the $3\tau$ cutoff as $z$ tends to zero. Near that
cutoff the coefficient of $T$ tends to zero and the coefficient of
$\mathbf e$ tends to one, uniformly on the clamped coordinate ranges.
Thus the computed value approaches default. If its radius estimate is
at least $3\tau$, it is already default. Coverage far from the curve
is not needed in either case.

For the outside prescription, use the special bits, the estimated mean
and outside squared radius in~\eqref{geo:outside}. Correct special
matching bits and good statistics make all its scalar and coordinate
arguments close to their exact values. The same protected denominators
and compact-range continuity apply, and $h_\mu$ is continuous. This
includes the special top, where the exact value is $-\mathbf e$.

Finally consider the choice between the inside and outside computations.
If it differs from the side of the true mean, a good
estimate forces $|\mu-1/2|$ to be small. In the interface slab all
endpoints remain far, and the only possible close piece is the always
included vertical. The inside computation is therefore valid for the
inside prescription extended to that slab. The outside computation is
valid for its extension using the clamped $h_\mu$. By
Lemma~\ref{geo:interface}, the two target values differ componentwise
by at most $6|\mu-1/2|$. Adding that error proves stability even with
the wrong side choice. These cases include exact mean, radius, and
minimum ties, without requiring a reliable hidden-real comparison at
zero gap.
\end{proof}

\subsection{Average accuracy and finite answer trees}
\label{geo:finite-schema}

\begin{proof}[Proof of Lemma~\ref{lem:geometric-procedures}]
\paragraph{Oracle accuracy with fixed references.}
Fix a desired $a>0$. Choose a component error $e>0$ and failure budget
$p>0$ so that
\begin{equation}\label{geo:accuracy-budget}
 4e^2+16B_1^2p<a^2.
\end{equation}
First choose the mesh, scalar tolerances, quantization width, and public
numerical precision for Lemma~\ref{geo:stability}, with slack for final
output rounding. Fix scalar sample counts that give its concentration
events a total failure probability below a chosen part of $p$.
The eight word roles, 32 piece roles, five possible length statistics,
mean, mesh-distance statistics, and outside-radius statistic now
determine a fixed number $N_{\rm call}$ of high-level completed calls. This count
includes all sampled bit requests and all matching final bit requests
over the possible relevant true references and flag alternatives.

Make validation and edge-flag errors small enough for their total budget.
For every remaining high-level bit call require average error smaller
than the allocated part of $p/N_{\rm call}$. Analyze each call on its possible
fixed reference before conditioning on acceptance or on the numerical
winner. Parts~\ref{geo:access-validation}--\ref{geo:access-statistics}
of the access contract and a union bound then charge the event of
\emph{any} required bit, decision, or statistical error by at most
$p$. Independence of these events is unnecessary. A fixed reference
does not have to be selected on every good trial; the unconditioned
analysis simply pays for its possible calls in advance. Validation
failures, including any acceptance with no proper reference, are
charged separately. The two special-word roles are charged in exactly
the same way when their supply is virtual.

On the complement of the charged failure event, component error is
at most $e$ by Lemma~\ref{geo:stability} and the planned final rounding.
On every outcome clamp the numerical answer to $[-B_1,B_1]$ and round
to a sufficiently fine fixed dyadic grid there. Detected invalid
numerical inputs may return any fixed bounded default. Undetected
incorrect choices are already included in the failure event. The true
field satisfies~\eqref{geo:bounded-field}, so the component error on
an arbitrary outcome is at most $2B_1$. Hence
\[
 4\E|\widehat v_I(X)-v_I(X)|^2
    \le4e^2+16B_1^2p<a^2.
\]
This proves~\eqref{geo:average-accuracy} unconditionally over the
output-coordinate and trial law. It does not assert correctness
conditioned on a rare selected coordinate or branch. The proof began
with arbitrary fixed admissible offsets, so it applies to each such
pair, and therefore also after averaging any chosen laws on them.

\paragraph{Finite oracle implementation.}
Here are the finite implementation and raw-address details. Fix a small
rational bin width $h>0$ and compare each physical scalar read to
thresholds $a+kh+\zeta$ on a fixed interval larger than $[-1,2]$,
where $\zeta$ is an offset modulo $h$. Clamp at the endpoints.
Representatives can depend just on the bin number, on a fixed rational
grid, and are within $O(h)$ of the input for every $\zeta$. This
error was included in the earlier tolerance choice. Comparisons whose
answers are forced by the allowed physical range can be public.
Seed bits use the common threshold $\xi$. Each offset may eventually
range over a fine finite rational grid. No numerical assertion above
requires a continuous offset distribution.

There are now only finitely many bits and bin indices in the geometric
computation. Public rational approximations to the fixed mesh sine and
cosine values, protected length normalizations, clamped square roots,
and all other numerical expressions can be tabulated to the chosen
precision. Bounded Boolean lookup implements the resulting answer
function. In particular neither an exact curve minimum nor an
unbounded hidden vertex name or transcript is processed by the local
Boolean computation. Public index arithmetic and any caller-specific
field or interpolation calculations are performed when constructing
the corresponding fixed branch and its wiring.

Wrong validation, flag, or type answers can invoke a reserved symbolic
routine on an unsuitable seed or select another reserved role. Wrong
length, radius, side, or minimum answers select a fixed mesh or lookup
case. None of these answers supplies a new word-coordinate address:
statistical coordinates remain the preallocated coordinates and final
data remain at $J$. Inactive or nonsensical roles may use public-zero
defaults, which require no query to a distinguished physical position.
Consequently direct scalar reads have the matching or prescribed
statistical marginals on every fixed hypothetical oracle branch. This
proves the structural assertion of part (iii), without inspecting the
physical realization of a supplied oracle answer. A constant number
of choices from polynomially enumerable oracle-level public spaces
has polynomial support; a constant product of rational probabilities
of polynomial bit complexity has polynomial bit complexity. Public
indices, oracle descriptions, and bounded answer tables retain the
same polynomial construction bound. This proves the oracle-level
enumerability assertion of part (iii).

\paragraph{Physical read load under the additional hypothesis.}
For part (iv), now assume the additional physical-read
hypothesis~\ref{geo:access-finite}. The supplied routines then have
its all-branch primitive bounds, even on unsuitable seeds.

Only at this point, after all requested accuracies and the repetitions
implementing them have been fixed, unfold the full lower-level answer
tree. Give unused possible nodes their prescribed public random choices
as well. It has a fixed finite number $H$ of potential primitive slots.
For a domain of size $M$, let $C_\ell/M$ bound the probability that
slot $\ell$ reads any specified position, for its fixed symbolic branch
under the unconditioned public law. A slot not using this domain has
$C_\ell=0$. Summing gives
\begin{equation}\label{geo:raw-load}
 \E_{I,\,\mathrm{trial}}
       [\text{number of potential reads of position }j]
 \le\frac{\sum_{\ell=1}^H C_\ell}{M}
 =\frac C M.
\end{equation}
This counts all wired possibilities rather than conditioning on the
executed path. Omitting an inactive occurrence only reduces its
unconditioned mass. In particular this argument does not require
accuracy on wrong branches. The high-level count $N_{\rm call}$ used for the
correctness union bound preceded amplification; the larger count $H$
is used for raw load and subsequent gate-failure accounting. It is
not fed back into the choice of the accuracy that determines $H$.

Refining offset grids changes the values of thresholds for a fixed
finite set of comparisons. It changes neither the bin width and bin
count, nor the mesh, nor the symbolic answer-tree slots, nor a
fixed branch's address primitive. Thus $H$ and $C$ remain uniform
under that refinement. Under the stated polynomial enumerability
premise for the physical routines, a constant number of their public
choices, also at unused tree nodes, again has polynomial support and
rational-probability bit complexity. Their public addresses and bounded
answer tables are constructed in polynomial time. This proves the
physical enumerability assertion of part (iv) and completes the proof.
\end{proof}

\begin{remark}[Late margins under physical realization]\label{geo:late-margins}
Assume the additional physical-read
hypothesis~\ref{geo:access-finite}. Its grid uniformity explains the
margin order for the fully expanded physical schema in part (iv).
Once the full number of physical comparison slots is fixed, a
uniform $K$-point threshold grid on an interval of length $L_0>0$
satisfies, for every fixed scalar $x$,
\[
 \Pr(|x-\xi|\le d)\le2/K+4d/L_0.
\]
For an offset grid modulo $h$, the probability of lying within $d$
of any bin boundary is at most $4/K+8d/h$. To see these bounds,
intersect the relevant interval of length $2d$ with the equally spaced
grid; each interval contains at most its length divided by the mesh
spacing plus two points. Modulo $h$ there are at most two intervals
after wraparound. These enlarged constants cover the endpoints.

Fix an output coordinate and all nonthreshold public choices on every
hypothetical branch. Each possible raw address is then fixed under
the address contract, regardless of the eventual branch. A union bound
over their fixed finite list, first choosing the offset grids fine
enough and then a positive rational margin small enough, makes the
probability of any small-gap comparison as small as prescribed,
uniformly in $X$. No search over offsets or enlargement of the
hidden-answer tree is involved. These margin losses are separate from
\eqref{geo:average-accuracy}; the latter already holds for each fixed
admissible offset. Occurrence-dependent perturbations within the margin
are allowed because the full list was charged before following any
branch. For an application using only the oracle contract, this
argument does not account for physical comparisons inside substituted
oracle gadgets. The outer construction uses its own direct-read and
functional-boundary margin analysis in
Section~\ref{sec:outer-composition}, together with the port theorem
and its separate physical substitution count.
\end{remark}

The geometric procedures are now proved independently of both circuit
constructions. Together with the port construction in
Section~\ref{sec:inner-gadget}, they establish the two interfaces used in
Section~\ref{sec:outer-composition}, completing the proof of
Theorem~\ref{thm:main}.
